\section{Conclusion}
\label{sec:conclusion}

In this paper, we investigate whether increasing the population size of an LLM ecosystem can prevent model collapse when models recursively train on shared synthetic data. We find that a larger population of homogeneous agents actually accelerates the degradation of generated text compared to a single-model baseline. However, introducing epistemic diversity through varied sampling temperatures provides a measurable mitigation against this collapse. Our findings demonstrate that a minimum viable population for an LLM ecosystem fundamentally requires active maintenance of diversity. Without varying epistemic viewpoints to counterbalance collective bias, an ecosystem will inevitably collapse, and larger homogeneous populations will simply reach that degraded state more efficiently. Future research should explore the effects of structural and architectural diversity over extended recursive generations to fully establish the conditions for a sustainable AI ecosystem.
